\section{Khái niệm Audit Trail/Audit Log và yêu cầu hệ thống}

\textbf{Audit Trail} (Audit Log) là tập hợp bản ghi có thứ tự thời gian ghi lại hoạt động hệ thống. Theo ISO/IEC 27002:2022 \cite{ISO27002}, Audit Log cần đảm bảo: (1) \textbf{Completeness} (Đầy đủ), (2) \textbf{Integrity} (Toàn vẹn), (3) \textbf{Non-repudiation} (Không thể phủ nhận), và (4) \textbf{Traceability} (Truy vết).

Trong CSDL quan hệ (\ac{oltp}), Audit Log ghi nhận giá trị trước và sau mỗi thao tác \ac{dml} (INSERT/UPDATE/DELETE). Ba đặc trưng kỹ thuật chính cần giải quyết:
\begin{itemize}
    \item \textbf{Write-heavy:} Dữ liệu sinh ra liên tục, đòi hỏi cơ chế quản lý vòng đời (retention, archiving) hiệu quả.
    \item \textbf{Latency overhead:} Cần giảm thiểu độ trễ (overhead $< 15\%$ \ac{tps}) để không ảnh hưởng nghiệp vụ.
    \item \textbf{Tính bất biến (WORM):} Log phải được bảo vệ khỏi sự thay đổi trái phép, kể cả từ DBA.
\end{itemize}

\section{Các giải pháp Audit Log hiện nay}

Các hệ thống cơ sở dữ liệu hiện đại có nhiều cách tiếp cận để xây dựng Audit Log:

\begin{itemize}
    \item \textbf{Logical Decoding (WAL-based / CDC):} Đọc Write-Ahead Log (WAL) và đẩy sự kiện ra ngoài (ví dụ: Debezium và Kafka \cite{Debezium}). \textit{Ưu điểm:} Throughput cực cao, không ảnh hưởng transaction. \textit{Nhược điểm:} Phức tạp, phụ thuộc hạ tầng ngoài (Kafka), không có SQL query trực tiếp và thường mất ngữ cảnh ứng dụng (application context).
    \item \textbf{pgAudit Extension:} Giải pháp chính thức ghi nhận statement-level \cite{pgAudit}. \textit{Ưu điểm:} Cài đặt nhanh, hỗ trợ \ac{pci}. \textit{Nhược điểm:} Ghi dưới dạng text file, không lưu trạng thái dữ liệu (OLD/NEW), không thể truy vấn có cấu trúc trực tiếp.
    \item \textbf{Database Trigger (Hstore/JSONB):} Ghi nhận thay đổi ngay trong CSDL thông qua Trigger \cite{Atkins2012}. Đây là hướng tiếp cận truyền thống nhưng được tối ưu mạnh mẽ nhờ kiểu dữ liệu JSONB. \textit{Ưu điểm:} Truy vấn trực tiếp bằng SQL, lưu đầy đủ ngữ cảnh ứng dụng, dễ triển khai. \textit{Nhược điểm:} Tăng overhead cho mỗi giao dịch.
\end{itemize}

\begin{center}
\captionof{table}{So sánh các giải pháp Audit Log}
\label{tab:solution-comparison}
\setlength{\tabcolsep}{4pt}
\footnotesize
\begin{tabularx}{\textwidth}{>{\raggedright\arraybackslash}p{3.5cm}>{\centering\arraybackslash}X>{\centering\arraybackslash}X>{\centering\arraybackslash}X}
\toprule
\textbf{Tiêu chí} & \textbf{WAL/ CDC} & \textbf{pgAudit} & \textbf{Trigger và JSONB} \\
\midrule
Structured Query & Không & Không & \textbf{Có} \\
OLD/NEW values & Có (raw) & Không & \textbf{Có (JSONB)} \\
App context & Không & Một phần & \textbf{Có} \\
Real-time streaming & Có & Không & Không \\
Hạ tầng ngoài & \textbf{Cao} & Không & Không \\
Phức tạp vận hành & Cao & Thấp & \textbf{Thấp} \\
\bottomrule
\end{tabularx}
\end{center}

\textbf{Lý do chọn Trigger, JSONB và Partitioning:} Đề tài ưu tiên truy vấn có cấu trúc bằng SQL, lưu trữ ngữ cảnh ứng dụng chính xác, và tích hợp cơ chế bảo mật (WORM, SHA-256) ngay bên trong PostgreSQL mà không phụ thuộc hạ tầng bên ngoài.

\section{Tại sao chọn PostgreSQL cho đề tài?}

PostgreSQL 16 cung cấp sẵn mọi tính năng (primitives) cần thiết để xây dựng giải pháp này mà không cần cài thêm extension ngoài:
\begin{itemize}
    \item \textbf{JSONB \& GIN Index} \cite{PG16JSONB, PG16GIN}: Cho phép lưu trữ và tìm kiếm dữ liệu phi cấu trúc hiệu suất cao.
    \item \textbf{Declarative Partitioning} \cite{PG16Partitioning}: Quản lý dữ liệu lớn bằng cách phân mảnh vật lý theo thời gian, tối ưu I/O và partition pruning.
    \item \textbf{SECURITY DEFINER \& Event Trigger}: Hỗ trợ phân quyền chặt chẽ và bắt các thay đổi cấu trúc dữ liệu (DDL).
    \item \textbf{pgcrypto \& dblink}: Cung cấp hàm băm SHA-256 cho chuỗi Hash và cơ chế autonomous transaction để cảnh báo bảo mật độc lập.
\end{itemize}
